% =====================================================================
% Style commun de l'affiche et de l'annexe (tikzposter, A0 portrait).
% Compiler avec XeLaTeX (ou LuaLaTeX) : les polices IBM Plex sont dans fonts/.
% =====================================================================
\usepackage{fontspec}
\usepackage{polyglossia}
\setdefaultlanguage{french}
\usepackage{graphicx}
\usepackage{booktabs}
\usepackage{array}
\usepackage{tabularx}
\usepackage{amsmath}
\usepackage{enumitem}
\usetikzlibrary{arrows.meta}

% ---------------------------------------------------------------- polices
\setmainfont{IBMPlexSans}[Path=fonts/, Extension=.ttf, UprightFont=*-Regular,
  BoldFont=*-SemiBold, ItalicFont=*-Italic]
\setsansfont{IBMPlexSans}[Path=fonts/, Extension=.ttf, UprightFont=*-Regular,
  BoldFont=*-SemiBold, ItalicFont=*-Italic]
\setmonofont{IBMPlexMono}[Path=fonts/, Extension=.ttf, UprightFont=*-Regular,
  BoldFont=*-SemiBold]
\newfontfamily\condensed{IBMPlexSansCondensed}[Path=fonts/, Extension=.ttf,
  UprightFont=*-Bold, BoldFont=*-Bold]
\newfontfamily\plexbold{IBMPlexSans}[Path=fonts/, Extension=.ttf, UprightFont=*-Bold]

% ---------------------------------------------------------------- couleurs
\definecolor{encre}{HTML}{16161A}
\definecolor{encre2}{HTML}{2C2C33}
\definecolor{gris}{HTML}{4A4A52}
\definecolor{filet}{HTML}{D9D6CE}
\definecolor{fond}{HTML}{FBFAF7}
\definecolor{carton}{HTML}{FFFFFF}
\definecolor{bordure}{HTML}{E4E2DC}
\definecolor{bleu}{HTML}{2A78D6}
\definecolor{orange}{HTML}{EB6834}
\definecolor{excit}{HTML}{3A3A42}
\definecolor{inhib}{HTML}{C2410C}
\definecolor{ambre}{HTML}{8A6D00}

% ---------------------------------------------------------------- theme tikzposter
\definecolorstyle{Affiche}{}{%
  \colorlet{backgroundcolor}{fond}
  \colorlet{framecolor}{fond}
  \colorlet{titlefgcolor}{encre}
  \colorlet{titlebgcolor}{fond}
  \colorlet{blocktitlebgcolor}{fond}
  \colorlet{blocktitlefgcolor}{encre}
  \colorlet{blockbodybgcolor}{fond}
  \colorlet{blockbodyfgcolor}{encre}
  \colorlet{innerblocktitlebgcolor}{fond}
  \colorlet{innerblocktitlefgcolor}{encre}
  \colorlet{innerblockbodybgcolor}{fond}
  \colorlet{innerblockbodyfgcolor}{encre}
  \colorlet{notefgcolor}{encre}
  \colorlet{notebgcolor}{carton}
  \colorlet{notefrcolor}{bordure}
}
\tikzposterlatexaffectionproofoff
\usecolorstyle{Affiche}
\usebackgroundstyle{Empty}
\usetitlestyle{Empty}
\useblockstyle{Minimal}
\defineblockstyle{Affiche}{titlewidthscale=1, bodywidthscale=1, titleleft,
  titleoffsetx=0pt, titleoffsety=0pt, bodyoffsetx=0pt, bodyoffsety=0pt,
  bodyverticalshift=0pt, roundedcorners=0, linewidth=0pt,
  titleinnersep=0mm, bodyinnersep=0mm}{}
\useblockstyle{Affiche}

% ---------------------------------------------------------------- commandes
% Titre de section numerote : \titre{01}{Contexte et objectifs}
\newcommand{\titre}[2]{{\condensed\fontsize{36pt}{40pt}\selectfont
  {\ttfamily\bfseries\color{bleu}#1}\hspace{0.5em}#2}\par\vspace{3mm}}
% Bandeau de partie, pleine largeur
\newcommand{\bandeau}[1]{\noindent{\ttfamily\bfseries\fontsize{20pt}{24pt}\selectfont\addfontfeature{LetterSpace=8}#1}%
  \hspace{1em}\leaders\hrule height 1.2mm\hfill\null}
% Legende de figure : \legende{Fig. 1 --- Titre.}{Texte}
\newcommand{\legende}[2]{\par\vspace{2mm}{\fontsize{17pt}{23pt}\selectfont\color{encre2}\textbf{#1} #2\par}\vspace{4mm}}
% Numero de piece cercle
\newcommand{\num}[1]{\tikz[baseline=(n.base)]\node[circle, draw=encre, line width=0.4mm, inner sep=0.6mm, font=\fontsize{13pt}{13pt}\selectfont\bfseries] (n) {#1};\ }
% Encadre blanc
\newcommand{\encadre}[1]{\noindent\colorbox{carton}{\fcolorbox{bordure}{carton}{\parbox{\dimexpr\linewidth-2\fboxsep-1.6pt-2\fboxrule\relax}{\raggedright #1}}}}
\newcommand{\code}[1]{\texttt{#1}}
\newcommand{\petit}{\fontsize{18pt}{24.5pt}\selectfont}
\newcommand{\corps}{\fontsize{19.5pt}{27pt}\selectfont}
\setlength{\fboxsep}{5mm}
\setlength{\fboxrule}{0.6mm}
\setlength{\parindent}{0pt}
\setlength{\parskip}{3mm}
\setlist{nosep, leftmargin=0pt, labelsep=0pt, itemsep=3mm}
\renewcommand{\arraystretch}{1.25}
\setlength{\heavyrulewidth}{0.6mm}
\setlength{\lightrulewidth}{0.3mm}

% Schemas TikZ
\tikzset{
  carte/.style={rectangle, rounded corners=1.6mm, line width=0.6mm, align=center, inner sep=1mm},
  pill/.style={rounded corners=8mm},
  etiq/.style={inner sep=0pt, text=gris},
  excitl/.style={draw=excit, line width=0.8mm},
  inhibit/.style={draw=inhib, line width=0.8mm},
  bande/.style={draw=bordure!80!gris, dash pattern=on 3mm off 2.5mm, line width=0.6mm, rounded corners=5mm, fill=bordure!35!fond},
}

% Motif raster du CPG (une barre bleue, une orange, periode 197 ms)
\newcommand{\raster}[1]{%
  \begin{tikzpicture}
    \foreach \i in {0,...,15} {
      \fill[bleu, rounded corners=0.4mm] (\i*5.2,0) rectangle ++(0.19,#1);
      \fill[orange, rounded corners=0.4mm] (\i*5.2+2.6,0) rectangle ++(0.19,#1);
    }
  \end{tikzpicture}}
